\begin{table}[htbp]
\centering
\caption{First Stage: March-2018 Reinterpretation and the SME-Composition Shift in Group-65 Winners}
\label{tab:first_stage}
\begin{adjustbox}{max width=\textwidth}
\begin{threeparttable}
\small
\begin{tabular}{lcccccc}
\toprule
 & (1) 6m, item & (2) 6m, item+PBU & (3) 12m, item & (4) 12m, item+PBU & (5) 18m, item & (6) 18m, item+PBU \\
\midrule
$g65 \times Pre$ & -0.0743*** & -0.0857*** & -0.1274*** & -0.1321*** & -0.1001*** & -0.1162*** \\
 & (0.0072) & (0.0072) & (0.0067) & (0.0066) & (0.0064) & (0.0063) \\
\midrule
Observations & 219,535 & 219,535 & 439,054 & 439,054 & 649,714 & 649,714 \\
Item FE & YES & YES & YES & YES & YES & YES \\
Month FE & YES & YES & YES & YES & YES & YES \\
PBU FE & NO & YES & NO & YES & NO & YES \\
\bottomrule
\end{tabular}
\begin{tablenotes}
\small
\item \textit{Notes:} Dependent variable is an indicator for whether
the winning firm was registered as ME or EPP. Sample is restricted to
completed items ($oc\_item\_status = 1$). $g65 \times Pre$ identifies
the Group-65 pre-treatment cell: a \emph{negative} coefficient means
Group-65 items had a lower SME-winner share \emph{before} the
reinterpretation than after, relative to the within-item average.
Equivalently, the March~2018 cutoff raised the SME-winner share in Group
65 by about 11--12 percentage points above what item and PBU fixed effects
alone would predict. Raw pre/post means: Group-65 pre = 0.1\%, post =
39.1\%; controls pre = 40.4\%, post = 67.0\% (controls also rise over
time, so the DiD is smaller than the within-G65 jump). This is the
\emph{first stage} of the identification strategy: the legal
reinterpretation measurably changes who wins Group-65 items. Standard
errors clustered at the item level in parentheses.
*** \textit{p}$<$0.01, ** \textit{p}$<$0.05, * \textit{p}$<$0.1.
\end{tablenotes}
\end{threeparttable}
\end{adjustbox}
\end{table}
